\section{Umklapp Scattering, Charge-Sector Leakage, and Finite Projection}
\label{sec:ch21-umklapp}

\subsection*{From a four-fermion process to a precise invariance question}

In the preceding sections the charge and spin currents were built from four
ordered fermion species, and the spin-singlet pair was found to have two
separate species-charge targets.  We now confront a related question for a
Hamiltonian rather than an order parameter: \emph{does a chosen finite
low-energy space remain invariant under a local interaction?}  An answer
requires more than the name ``Umklapp.''  It requires an ordered CAR word,
an explicit charge shift, and a state on which the operator acts nontrivially
\emph{outside} the chosen space.  Pauli exclusion can make a seemingly
outward process vanish; the opposite process may remain nonzero while staying
inside the permitted charge box.

The discussion is entirely on a finite lattice and uses the occupation
representation from Chapter~\ref{ch:car_fock}.  We prove exact selection rules,
self-adjointness, and a conditional failure of invariance.  These statements
do not imply a thermodynamic charge gap or establish continuum renormalization
group equations.  For physical motivation and conventional bosonized models,
see, for example, \cite{giamarchi2004,gogolin1998}; our claims below are
finite algebraic statements with their hypotheses displayed.

As in Sections~\ref{sec:ch21-spin-charge} and
\ref{sec:ch21-singlet-pair}, choose \(h\in\mathbb N_{>0}\) and set
\begin{equation}
  L=2h,\qquad \mathcal B_L=(-h,h]\cap\mathbb Z
       =\{1-h,\ldots,h\},\qquad
  \Lambda_L=\mathbb Z/L\mathbb Z,\qquad \zeta=e^{2\pi i/L}.
  \label{eq:ch21-u-band}
\end{equation}
The four species have the fixed occupation ordering
\begin{equation}
 L\uparrow<L\downarrow<R\uparrow<R\downarrow,
 \quad\text{with increasing momentum inside each species}.
 \label{eq:ch21-u-order}
\end{equation}
For \(\alpha\in\mathcal C=\{L,R\}\times\{\uparrow,\downarrow\}\),
let \(c_{\alpha,k}\) satisfy the finite canonical anticommutation relations
(CAR).  We use periodic position fields in the Fourier convention
\begin{equation}
 c_{\alpha,x}=\frac1{\sqrt L}\sum_{k\in\mathcal B_L}
                \zeta^{kx}c_{\alpha,k},\qquad
 c_{\alpha,x}^{\dagger}=\frac1{\sqrt L}\sum_{k\in\mathcal B_L}
                \zeta^{-kx}c_{\alpha,k}^{\dagger}.
 \label{eq:ch21-u-site-fields}
\end{equation}
The positive-mode cutoff, if density excitations are considered separately,
is denoted \(M\) (with \(1\leq m\leq M\)); no density-mode or
current-CCR cutoff enters the CAR calculation in this section.

\subsection*{The ordered Umklapp word and its Hermitian Hamiltonian}

\begin{definition}[Finite-site Umklapp interaction]
\label{def:ch21-u-word}
Define the ordered product
\begin{equation}
  O_U(x)=c_{L\uparrow,x}^{\dagger}
         c_{L\downarrow,x}^{\dagger}
         c_{R\downarrow,x}c_{R\uparrow,x},
       \qquad x\in\Lambda_L.
  \label{eq:ch21-u-word}
\end{equation}
For a \emph{real} coupling \(g_U\in\mathbb R\), set
\begin{equation}
 H_U=\frac{g_U}{L}\sum_{x\in\Lambda_L}
           \bigl(O_U(x)+O_U(x)^\dagger\bigr).
 \label{eq:ch21-u-hamiltonian}
\end{equation}
Products act from right to left: the right-up and right-down particles
are removed, in that order, before the left-down and left-up particles are
created.  The factor \(1/L\) in \eqref{eq:ch21-u-hamiltonian} is fixed
throughout, including the finite coefficient examples below.
\end{definition}

In the finite Fock representation these are bounded operators on the
same Hilbert space \(\mathcal H_L\), so no operator-domain problem arises.
The adjoint reverses \emph{the order} of the four factors:
\begin{equation}
 O_U(x)^\dagger
  =c_{R\uparrow,x}^\dagger c_{R\downarrow,x}^\dagger
    c_{L\downarrow,x}c_{L\uparrow,x}.
 \label{eq:ch21-u-adjoint}
\end{equation}
It is not permissible to reorder these odd operators without keeping
the resulting CAR signs.

\begin{proposition}[Self-adjointness]
\label{prop:ch21-u-selfadjoint}
For every \(h>0\) and real \(g_U\), the operator \(H_U\) in
\eqref{eq:ch21-u-hamiltonian} satisfies \(H_U^\dagger=H_U\).
\end{proposition}
\begin{proof}
Write \(A=\sum_xO_U(x)\).  Since the sum is finite,
\(A^\dagger=\sum_xO_U(x)^\dagger\); hence
\[
 H_U^\dagger
 =\overline{\frac{g_U}{L}}(A^\dagger+A)
 =\frac{g_U}{L}(A+A^\dagger)=H_U.
\]
The middle equality uses the reality of \(g_U\).  No commutation
between \(A\) and \(A^\dagger\) was assumed.
\end{proof}

An Umklapp word is a possible lattice process rather than a statement of
its nonzero action on all states.  Its momentum-space expansion follows
exactly from \eqref{eq:ch21-u-site-fields}:
\begin{align}
 O_U(x)=\frac1{L^2}\!
 \sum_{k_1,k_2,k_3,k_4\in\mathcal B_L}
 &\zeta^{(-k_1-k_2+k_3+k_4)x}\,
 c_{L\uparrow,k_1}^\dagger c_{L\downarrow,k_2}^\dagger
 c_{R\downarrow,k_3}c_{R\uparrow,k_4}.
 \label{eq:ch21-u-fourier-expansion}
\end{align}
Summation over \(x\) retains precisely those ordered words satisfying
\(k_1+k_2-k_3-k_4\equiv0\pmod L\).  Thus momentum is conserved
\emph{modulo} the reciprocal-lattice period, rather than as an
unrestricted equality between the chosen integer representatives.
This distinction will be essential for the explicit witness.

\subsection*{Charge shifts and the surviving conservation laws}

The relative species number operators are
\begin{equation}
 \widehat N_\alpha=\sum_{k\in\mathcal B_L}
     c_{\alpha,k}^\dagger c_{\alpha,k}-hI,
 \qquad \alpha\in\mathcal C.
 \label{eq:ch21-u-relative-charge}
\end{equation}
They have integer eigenvalues \(-h,\ldots,h\) on the finite Fock space.
Let \(e_\alpha\) denote a basis vector in the four-dimensional
species-charge lattice.  With the order \eqref{eq:ch21-u-order} we define
\begin{equation}
 \Delta\mathbf N
  =e_{L\uparrow}+e_{L\downarrow}
     -e_{R\uparrow}-e_{R\downarrow}
  =(+1,+1,-1,-1).
 \label{eq:ch21-u-charge-vector}
\end{equation}
Every \emph{species} shift has magnitude one, whereas each chiral
\emph{branch total} changes by two.

\begin{lemma}[One-particle charge commutators]
\label{lem:ch21-u-fundamental-commutators}
For any species \(\gamma,\beta\), momentum \(q\in\mathcal B_L\), and
site \(x\),
\begin{align}
 [\widehat N_\gamma,c_{\beta,q}]
   &=-\delta_{\gamma\beta}c_{\beta,q},&
 [\widehat N_\gamma,c_{\beta,q}^\dagger]
   &=\delta_{\gamma\beta}c_{\beta,q}^\dagger,
 \label{eq:ch21-u-elementary-commutators}\\
 [\widehat N_\gamma,c_{\beta,x}]
   &=-\delta_{\gamma\beta}c_{\beta,x},&
 [\widehat N_\gamma,c_{\beta,x}^\dagger]
   &=\delta_{\gamma\beta}c_{\beta,x}^\dagger.
 \label{eq:ch21-u-site-commutators}
\end{align}
\end{lemma}
\begin{proof}
The CAR imply, for individual modes \(i,j\),
\[
 [c_i^\dagger c_i,c_j]
   =c_i^\dagger\{c_i,c_j\}-\{c_i^\dagger,c_j\}c_i
   =-\delta_{ij}c_j,
\]
and similarly \([c_i^\dagger c_i,c_j^\dagger]=\delta_{ij}c_j^\dagger\).
Summing over the momenta of species \(\gamma\) proves
\eqref{eq:ch21-u-elementary-commutators}; the scalar sea subtraction
commutes with everything.  Fourier linearity gives
\eqref{eq:ch21-u-site-commutators}.
\end{proof}

\begin{theorem}[Exact four-species Umklapp selection rule]
\label{thm:ch21-u-charge-shift}
For every \(x\) and species \(\gamma\),
\begin{align}
 [\widehat N_\gamma,O_U(x)]
   &=\Delta N_\gamma\,O_U(x),
 \label{eq:ch21-u-commutator}\\
 [\widehat N_\gamma,O_U(x)^\dagger]
   &=-\Delta N_\gamma\,O_U(x)^\dagger.
 \label{eq:ch21-u-adjoint-commutator}
\end{align}
Consequently, if \(\mathcal H_{\mathbf N}\) is the simultaneous
relative-charge eigenspace of \(\mathbf N\in\mathbb Z^{\mathcal C}\), then
\begin{equation}
 O_U(x)\mathcal H_{\mathbf N}\subseteq
     \mathcal H_{\mathbf N+\Delta\mathbf N},\qquad
 O_U(x)^\dagger\mathcal H_{\mathbf N}\subseteq
     \mathcal H_{\mathbf N-\Delta\mathbf N}.
 \label{eq:ch21-u-charge-targets}
\end{equation}
If a listed charge is outside the allowed finite range, its eigenspace
is understood to be zero.  The total relative particle number and the
spin imbalance satisfy
\begin{equation}
 \left[\sum_{\alpha\in\mathcal C}\widehat N_\alpha,H_U\right]=0,
 \qquad
 [\widehat S,H_U]=0,
 \quad \widehat S=\sum_{\eta=L,R}
        (\widehat N_{\eta\uparrow}-\widehat N_{\eta\downarrow}).
 \label{eq:ch21-u-conserved-totals}
\end{equation}
By contrast, \(H_U\) need not conserve each species charge or the
separate left/right branch charge.
\end{theorem}
\begin{proof}
For arbitrary operators, the commutator is a derivation:
\([D,AB]=[D,A]B+A[D,B]\).  Expanding the ordered word gives
\begin{align*}
 [\widehat N_\gamma,O_U(x)]={}&
   [\widehat N_\gamma,c_{L\uparrow,x}^\dagger]
           c_{L\downarrow,x}^\dagger c_{R\downarrow,x}c_{R\uparrow,x}
 +c_{L\uparrow,x}^\dagger
   [\widehat N_\gamma,c_{L\downarrow,x}^\dagger]
           c_{R\downarrow,x}c_{R\uparrow,x}\\
 &+c_{L\uparrow,x}^\dagger c_{L\downarrow,x}^\dagger
   [\widehat N_\gamma,c_{R\downarrow,x}]c_{R\uparrow,x}
 +c_{L\uparrow,x}^\dagger c_{L\downarrow,x}^\dagger
   c_{R\downarrow,x}[\widehat N_\gamma,c_{R\uparrow,x}].
\end{align*}
Lemma~\ref{lem:ch21-u-fundamental-commutators} reduces the coefficient
to
\(\delta_{\gamma,L\uparrow}+\delta_{\gamma,L\downarrow}
-\delta_{\gamma,R\downarrow}-\delta_{\gamma,R\uparrow}\), which is
exactly \(\Delta N_\gamma\).  Taking adjoints of the operator
commutator, using self-adjointness of \(\widehat N_\gamma\), gives
\eqref{eq:ch21-u-adjoint-commutator}.  For
\(\widehat N_\gamma\psi=N_\gamma\psi\),
\[
 \widehat N_\gamma O_U(x)\psi
  =(N_\gamma+\Delta N_\gamma)O_U(x)\psi,
\]
proving the first inclusion in \eqref{eq:ch21-u-charge-targets};
the second is identical with the opposite shift.  Finally the
coefficients of the total-number commutator sum to \(1+1-1-1=0\),
while those of \(\widehat S\) sum to \((1-1)+(-1+1)=0\).
The identities persist under the finite site sum and the adjoint.
\end{proof}

Writing \(\widehat Q_c^\eta=\widehat N_{\eta\uparrow}
 +\widehat N_{\eta\downarrow}\) and
\(\widehat Q_s^\eta=\widehat N_{\eta\uparrow}
 -\widehat N_{\eta\downarrow}\) as in \S\,21.1, the same computation
immediately gives
\begin{equation}
 [\widehat Q_c^L,O_U]=2O_U,\qquad
 [\widehat Q_c^R,O_U]=-2O_U,\qquad
 [\widehat Q_s^L,O_U]=[\widehat Q_s^R,O_U]=0.
 \label{eq:ch21-u-branch-shifts}
\end{equation}
Here \(O_U\) denotes either the fixed-site word or its outward site sum.
These equations describe an interbranch transfer of two particles,
not the creation or destruction of total electric charge.

\begin{corollary}[Vanishing same-sector compression]
\label{cor:ch21-u-sector-compression}
For any admissible charge vector \(\mathbf N\) and any two nonnegative
energy cutoffs \(K,K'\),
\begin{equation}
 P_{\mathbf N,K}\,H_U\,P_{\mathbf N,K'}=0.
 \label{eq:ch21-u-zero-compression}
\end{equation}
\end{corollary}
\begin{proof}
Both projected source and projected target are contained in
\(\mathcal H_{\mathbf N}\).  Theorem~\ref{thm:ch21-u-charge-shift}
places the two summands of \(H_U\) in
\(\mathcal H_{\mathbf N+\Delta\mathbf N}\) and
\(\mathcal H_{\mathbf N-\Delta\mathbf N}\), respectively.  Neither
charge vector equals \(\mathbf N\), and joint eigenspaces with
different charge labels are orthogonal.  The left projection therefore
kills both contributions, independently of \(K,K'\).
\end{proof}

\subsection*{Charge boxes, projections, and a conditional leakage theorem}

A single-charge-sector projection is not the same as a charge
\emph{box}, which admits several sectors but limits their extent.
To define the latter without using an informal inequality between
vectors, let \(S\subseteq\mathcal C\times\mathcal B_L\) be an
occupation configuration.  Its relative charge vector is
\(\mathbf N(S)\), and its nonnegative integer excitation energy
\(\mathcal E(S)\) is measured above the ground configuration of the
\emph{same} charge vector, as in Chapter~\ref{ch:vacuum_budget}.

\begin{definition}[Per-species charge box and budget]
\label{def:ch21-u-charge-box}
Given bounds \(B_\alpha\in\mathbb N\) and an energy budget
\(K\in\mathbb N\), let
\begin{equation}
 \mathcal B_{K,\mathbf B}
   =\operatorname{span}\Bigl\{|S\rangle:\;
       |N_\alpha(S)|\le B_\alpha\ \text{for all }\alpha,
       \quad\mathcal E(S)\le K\Bigr\}
       \subseteq\mathcal H_L,
 \label{eq:ch21-u-box}
\end{equation}
and let \(P_{K,\mathbf B}\) be its orthogonal coordinate projection.
The charge bounds and the energy cutoff play logically distinct roles.
A box may be proper even when the energy cutoff is large enough to
include every state whose charge vector satisfies the bounds.
\end{definition}

\begin{theorem}[Conditional outward leakage]
\label{thm:ch21-u-leakage}
Assume \(h\ge1\), \(g_U\in\mathbb R\setminus\{0\}\), and
\(\psi\in\mathcal B_{K,\mathbf B}\) has definite charge
\(\mathbf N\), i.e.\ \(\widehat N_\alpha\psi=N_\alpha\psi\)
for all \(\alpha\).  Suppose
\begin{equation}
  \sum_{x\in\Lambda_L}O_U(x)\psi\ne0,
  \qquad
  \exists\alpha\in\mathcal C:\quad
      |N_\alpha+\Delta N_\alpha|>B_\alpha.
 \label{eq:ch21-u-leakage-hypotheses}
\end{equation}
Then \(H_U\psi\notin\mathcal B_{K,\mathbf B}\), and the ambient
Hamiltonian does not preserve this budget.  In addition,
\begin{equation}
 [P_{K,\mathbf B},H_U]\ne0.
 \label{eq:ch21-u-projector-noncommute}
\end{equation}
An analogous conclusion holds for a nonzero adjoint (inward) site sum
whose target \(\mathbf N-\Delta\mathbf N\) violates a charge bound.
\end{theorem}
\begin{proof}
For any \(\mathbf T\), denote the orthogonal full-charge-sector
projection by \(\Pi_{\mathbf T}\).  Distinct sectors are orthogonal
because the \(\widehat N_\alpha\) are commuting self-adjoint matrices.
The two shifts in \eqref{eq:ch21-u-charge-targets} differ, hence
\begin{equation}
 \Pi_{\mathbf N+\Delta\mathbf N}H_U\psi
   =\frac{g_U}{L}\sum_xO_U(x)\psi\ne0.
 \label{eq:ch21-u-projected-outward}
\end{equation}
The adjoint contribution has charge \(\mathbf N-\Delta\mathbf N\)
and cannot cancel the displayed component.  By the second hypothesis,
\(\mathcal B_{K,\mathbf B}\) has no vector in the charge sector
\(\mathbf N+\Delta\mathbf N\).  Equation
\eqref{eq:ch21-u-projected-outward} therefore exhibits a nonzero
component of \((I-P_{K,\mathbf B})H_U\psi\).
Since \(P_{K,\mathbf B}\psi=\psi\), we also have
\[
 [P_{K,\mathbf B},H_U]\psi
  =P_{K,\mathbf B}H_U\psi-H_UP_{K,\mathbf B}\psi
  =-(I-P_{K,\mathbf B})H_U\psi\ne0.
\]
The proof for the opposite shift exchanges \(O_U\) with \(O_U^\dagger\).
\end{proof}

The same elementary argument detects \emph{energy} leakage: it suffices
to exhibit an occupation ket \(|T\rangle\) excluded by the energy
cutoff for which \(\langle T,H_U\psi\rangle\ne0\).
The charge-box hypothesis in Theorem~\ref{thm:ch21-u-leakage} is a
convenient sufficient condition, not a necessary condition for
non-invariance.  Conversely, \(H_U\psi\ne0\) by itself is
\emph{not} a leakage certificate.  The latter distinction will be
visible in the second finite example.

\subsection*{An explicit outward witness: the sign and the Fourier coefficient}

The next calculation checks every part of the leakage hypothesis with
a concrete occupation configuration.  It also illustrates why a
fermionic sign cannot be inferred from charge conservation alone.

\begin{theorem}[Zero-charge sea leaves every zero-charge box]
\label{thm:ch21-u-zero-box-witness}
For every \(h\ge1\), write \(L=2h\).  In each species fill the sea
momenta \(\{1-h,\ldots,0\}\), and let \(\Omega\) be the corresponding
normalized, ordered occupation ket.  Define \(T\) by removing
momentum \(1-h\) from \(R\uparrow\) and \(R\downarrow\), and
inserting momentum \(1\) in \(L\uparrow\) and \(L\downarrow\).
Then
\begin{equation}
 \langle T,H_U\Omega\rangle=-\frac{g_U}{L^2}.
 \label{eq:ch21-u-witness-coefficient}
\end{equation}
If \(g_U\ne0\), the zero-charge sea \(\Omega\) belongs to
\(\mathcal B_{0,\mathbf0}\), but
\(H_U\Omega\notin\mathcal B_{K,\mathbf0}\) for \emph{every}
\(K\ge0\).  Thus \(H_U\) does not preserve the zero-charge budget
and does not commute with its projector.
\end{theorem}
\begin{proof}
Let \(S_0\) be the set of \(4h\) occupied species-momentum modes
constituting \(\Omega\).  The target configuration is
\[
 T=\bigl(S_0\setminus\{(R\uparrow,1-h),(R\downarrow,1-h)\}\bigr)
     \cup\{(L\uparrow,1),(L\downarrow,1)\}.
\]
For \(h\ge1\), the removed momenta lie in the sea and the inserted
momentum \(1\) lies in the band but outside the sea.  Therefore no
Pauli factor kills this particular word.  The relative-charge vector
of \(T\) is \(\Delta\mathbf N=(1,1,-1,-1)\), whereas
\(\mathbf N(\Omega)=\mathbf0\).

For an occupied mode \(i\), annihilation contributes the sign
\((-1)^{\#\{j\in S:j<i\}}\), and creation at an unoccupied mode
contributes the corresponding sign of the updated configuration.
The ordered word \eqref{eq:ch21-u-word} acts from right to left.
The number of occupied modes \emph{preceding} each operation, after
all previous operations have been made, is
\begin{center}
\begin{tabular}{@{}lll@{}}
\toprule
Operation & Preceding occupied modes & Sign \\
\midrule
\(c_{R\uparrow,1-h}\) & \(2h\) & \((-1)^{2h}\) \\
\(c_{R\downarrow,1-h}\) & \(3h-1\) & \((-1)^{3h-1}\) \\
\(c_{L\downarrow,1}^\dagger\) & \(2h\) & \((-1)^{2h}\) \\
\(c_{L\uparrow,1}^\dagger\) & \(h\) & \((-1)^h\) \\
\bottomrule
\end{tabular}
\end{center}
For example, the first removal has the two filled left species before
it; after that removal the right-down species has \(2h+(h-1)\)
occupied modes before its lowest sea momentum.  The insertion in
left-down sees all \(h\) left-up and its own \(h\) filled sea modes,
while the final insertion in left-up sees its original \(h\) sea modes.
The exponent is \(2h+(3h-1)+2h+h=8h-1\), an odd integer.
Thus the selected momentum word produces \(-|T\rangle\).

In \eqref{eq:ch21-u-fourier-expansion}, the four momenta are
\(k_1=k_2=1\) and \(k_3=k_4=1-h\).  Their Fourier phase is
\[
 \zeta^{(-1-1+(1-h)+(1-h))x}
    =\zeta^{-2hx}=\zeta^{-Lx}=1
       \quad(x\in\Lambda_L).
\]
Each site word has scalar Fourier normalization \(L^{-2}\).
The four changed occupations uniquely fix the four participating
momenta: another term in the momentum sum cannot produce \(|T\rangle\)
from \(|\Omega\rangle\).  The adjoint word changes the charge by
\(-\Delta\mathbf N\), hence also cannot contribute to this target.
The \(L\) equal site contributions multiplied by the Hamiltonian
prefactor \(g_U/L\) therefore give
\[
 \langle T,H_U\Omega\rangle
  =\frac{g_U}{L}\,L\,\frac{-1}{L^2}
  =-\frac{g_U}{L^2}.
\]
This is nonzero when \(g_U\ne0\), and \(|T\rangle\) is orthogonal
to every zero-species-charge box.  Theorem~\ref{thm:ch21-u-leakage}
then proves the final assertions.  No frozen-margin assumption for
finite density currents has been used.
\end{proof}

\begin{example}[Four-site value]
\label{ex:ch21-u-four-site}
At \(h=2\), \(L=4\), \(\zeta=i\), and \(g_U=1\),
the sea fills \(\{-1,0\}\) in every species.  The witness target
fills \(\{-1,0,1\}\) in each left species and \(\{0\}\) in
each right species.  Equation~\eqref{eq:ch21-u-witness-coefficient}
gives the exact matrix element
\[
 \langle T,H_U\Omega\rangle=-\frac1{16}.
\]
The charge vector \((1,1,-1,-1)\) already suffices to exclude this
ket from the zero-charge box, independently of the excitation budget.
\end{example}

\subsection*{Why a nonzero action need not leak: a Pauli-blocked counterexample}

\begin{example}[Only the inward channel survives]
\label{ex:ch21-u-inward-only}
Take \(h=2\), \(L=4\), and \(g_U\ne0\).  Choose a normalized occupation
ket \(|S\rangle\) with both left species \emph{fully} occupied,
\(\{-1,0,1,2\}\), and both right species occupied only at \(\{-1\}\).
Its relative charges are
\[
  \mathbf N(S)=(2,2,-1,-1).
\]
Use the proper charge box with bounds
\(\mathbf B=(2,2,1,1)\) in the species order
\((L\uparrow,L\downarrow,R\uparrow,R\downarrow)\), and the energy
cutoff \(K=16\).  The unprojected outward sum annihilates \(|S\rangle\):
\[
 O_U(x)|S\rangle=0\qquad\text{for every }x,
\]
because it would attempt to create a fermion in an already full
left species.  On the other hand, the adjoint sum is nonzero and
shifts the charges to \((1,1,0,0)\), which obey the same bounds.

To exhibit a coefficient, let \(|T'\rangle\) be the occupation ket with
left species \(\{0,1,2\}\), right-up \(\{-1,2\}\), and
right-down \(\{-1,0\}\).  The adjoint word removes \(-1\) from
each left species and inserts \(2\) in right-up and \(0\) in
right-down.  Its right-to-left CAR preceding counts are
\(0,3,8,7\), with an even sum \(18\), so the sign is positive.
The Fourier phase is \(\zeta^{(-1-1-2-0)x}=\zeta^{-4x}=1\).
Exactly as in the preceding proof,
\begin{equation}
 \langle T',H_U S\rangle=\frac{g_U}{16}\ne0.
 \label{eq:ch21-u-inward-coefficient}
\end{equation}
For this four-site band, the excitation energy above a fixed-number
sector ground is at most \(n(4-n)\le4\) for a species with \(n\)
occupied modes; summing four species gives \(\mathcal E\le16\).
Thus \emph{every} ket generated from \(|S\rangle\) by the adjoint
Umklapp term belongs to \(\mathcal B_{16,\mathbf B}\).  We have
\(H_U|S\rangle\ne0\) but \((I-P_{16,\mathbf B})H_U|S\rangle=0\).
This proves only that \emph{this state} does not leak; it does not
establish invariance of the whole box under \(H_U\).
\end{example}

\subsection*{What finite leakage does, and does not, establish}

The exact conclusions of this section can now be separated.  The CAR
calculation proves a four-species selection rule and identifies precisely
which charge sectors an Umklapp interaction couples.  Same-sector
first-order compression vanishes, but a charge-bounded or energy-bounded
subspace can fail to be invariant, as witnessed by a nonzero
\emph{outside} matrix element.  The finite compression
\(P_{K,\mathbf B}H_UP_{K,\mathbf B}\) remains a legitimate Hermitian
matrix on the retained carrier; it is simply not the restriction of an
invariant ambient dynamics when the leakage is nonzero.

The finite energy-shell machinery developed in Section~\ref{sec:ch21-finite-effective}
can subsequently be used to organize perturbative block reduction. Its
quadratic Schrieffer--Wolff/Baker--Campbell--Hausdorff identities concern
stated projectors and polynomial coefficients through a fixed order; they do
not automatically provide an analytic conjugation, error estimate, or
RG semigroup.  In particular, removing a shell of \emph{total excitation
energy} is not identical to removing a particular Fourier momentum
mode.  The shell-compression formalism must be instantiated with its
projector and block-diagonality hypotheses before being applied to
\(H_U\).

Finally, a Mott gap would require a spectral statement for a specified
many-body Hamiltonian, and a physically meaningful thermodynamic claim
would require a controlled sequence of systems as \(L\to\infty\),
with a specified density, coupling scaling, and limiting observable.
Likewise the Kosterlitz--Thouless flow equations belong to a
separately constructed continuum effective theory; the presence of a
nonzero finite matrix element \(-g_U/L^2\) neither derives their beta
functions nor proves gap formation.  Any proposed use of bosonized
charge vertices also requires care with the distinction between a
finite-CAR operator, an abstract oscillator representation, and a
coefficientwise formal series.  Chapter~21 remains unimplemented in
Lean at the time of this draft: prior Core support for CAR, charge
sectors and finite projections is not a verification of the new
Umklapp theorems.
